\section{Workflow Architecture Của Hệ Thống CD}

Phần này mô tả logic vận hành của đồ án ở mức workflow. Mục tiêu chính không chỉ là
``deploy được'', mà là mọi thay đổi runtime đều có provenance: commit nguồn trong app
repo, image tag trên Docker Hub, commit GitOps trong CD repo, và trạng thái reconcile
trong ArgoCD.

\subsection{Sơ Đồ Source-To-Runtime}

\begin{figure}[H]
\centering
\fbox{\begin{minipage}{0.94\textwidth}
\centering
\textbf{Developer / GitHub}\\
\path{tzin1401/yas}: branch, main, hoặc tag \texttt{vX.Y.Z}\\[0.6em]
$\Downarrow$\\[-0.1em]
\textbf{Jenkins App Pipeline}\\
Detect changed modules $\rightarrow$ Gitleaks/Test/Coverage/SonarQube $\rightarrow$ Maven build $\rightarrow$ Docker image\\[0.6em]
$\Downarrow$\\[-0.1em]
\textbf{Docker Hub}\\
\path{docker.io/<user>/yas-<service>:<commit-sha|main|latest|vX.Y.Z>}\\[0.6em]
$\Downarrow$\\[-0.1em]
\textbf{GitOps Commit}\\
\path{emanhthangngot/yas-cd}: update \path{overlays/<env>/kustomization.yaml}\\[0.6em]
$\Downarrow$\\[-0.1em]
\textbf{ArgoCD}\\
\texttt{yas-dev} auto-sync, \texttt{yas-staging} manual-sync, \texttt{yas-developer} on-demand\\[0.6em]
$\Downarrow$\\[-0.1em]
\textbf{K3s Runtime}\\
Namespaces \texttt{dev}, \texttt{staging}, \texttt{developer}, platform services, Istio mesh, NodePort 30846
\end{minipage}}
\caption{Workflow source-to-runtime của hệ thống CD}
\end{figure}

\subsection{Ranh Giới Trách Nhiệm Giữa Hai Repo}

\begin{table}[H]
\centering
\caption{Boundary giữa app repo, CD repo và runtime}
\begin{tabular}{|p{3.3cm}|p{5.2cm}|p{5.4cm}|}
\hline
\textbf{Lớp} & \textbf{Chịu trách nhiệm} & \textbf{Không chịu trách nhiệm} \\
\hline
App/CI repo & Source code YAS, Jenkinsfile chính, CI gate, Docker build/push, \path{services.yaml}, app instrumentation. & Không lưu desired state cuối cùng của cluster; không apply trực tiếp vào namespace ArgoCD quản lý. \\
\hline
CD/GitOps repo & Kustomize base/overlays, chart snapshot, ArgoCD apps, Jenkins job GitOps nhẹ, sampledata seed job, evidence/runbook/report. & Không build source code Java/Next.js chính; không thay thế CI gate của app repo. \\
\hline
Runtime K3s & Chạy workload, platform services, ingress, mesh, observability, ArgoCD reconcile. & Không là nơi chỉnh image bằng tay; trạng thái mong muốn phải quay về Git. \\
\hline
\end{tabular}
\end{table}

\subsection{Luồng Quyết Định Theo Trigger}

\begin{table}[H]
\centering
\caption{Trigger, tag image và chính sách sync}
\begin{tabular}{|p{2.7cm}|p{3.7cm}|p{4.2cm}|p{3.4cm}|}
\hline
\textbf{Trigger} & \textbf{Image tag} & \textbf{GitOps action} & \textbf{ArgoCD policy} \\
\hline
Feature branch & Short commit SHA 12 ký tự. & Không update dev/staging; developer preview dùng job riêng để chọn service/branch. & \texttt{yas-developer} auto-sync khi preview active. \\
\hline
\texttt{main} & Commit SHA, \texttt{main}, \texttt{latest}. & Update \path{overlays/dev} bằng commit SHA để tạo manifest diff thật. & \texttt{yas-dev} auto-sync/self-heal. \\
\hline
\texttt{vX.Y.Z} & Promote image hiện có sang release tag. & Update \path{overlays/staging} bằng immutable tag. & \texttt{yas-staging} manual-sync để operator duyệt. \\
\hline
Teardown developer & Không build image mới. & Reset \path{overlays/developer} về dormant/default tags. & \texttt{yas-developer} reconcile về \texttt{0/0}. \\
\hline
\end{tabular}
\end{table}

\subsection{Rollback Và Audit Trail}

Rollback trong mô hình này không dùng \texttt{kubectl set image}. Nhóm có thể rollback bằng
cách revert GitOps commit hoặc chỉnh lại tag trong overlay rồi để ArgoCD reconcile. Chuỗi
audit cần đọc theo thứ tự:

\begin{enumerate}
    \item App commit hoặc release tag trong \path{tzin1401/yas}.
    \item Jenkins console log chứng minh CI gate/build/promote.
    \item Docker Hub tag tồn tại cho service tương ứng.
    \item GitOps commit trong \path{emanhthangngot/yas-cd}.
    \item ArgoCD sync status và Kubernetes rollout/runtime smoke.
\end{enumerate}

Đây là điểm khác biệt quan trọng so với cách deploy thủ công: trạng thái runtime có thể giải
thích ngược về Git và image tag, không phụ thuộc thao tác trên cluster.
